% latexmk -pdf -pdflatex="lualatex %O %S" enos-slides.tex

% ************************************************************************
% Personal data and user ad-hoc commands
% ************************************************************************
\newcommand{\myName}{%
  Ronan-Alexandre Cherrueau, Adrien Lebre,\newline Dimitri Pertin,
  \underline{Anthony Simonet} and Matthieu
  Simonin\xspace} % 2
\newcommand{\myDepartment}{Inria, France} % 8
\newcommand{\myTime}{\date\xspace} % 9
\newcommand{\myTitle}{Enos: a Holistic Framework for Conducting
  Scientific Evaluations of OpenStack} % 11
\newcommand{\myVersion}{version 0.1\xspace}
\date{May 14th, 2017}
 \titlegraphic{%
	\includegraphics[height=1.0cm]{figs/logo/inria.pdf}\hfill%
	\includegraphics[height=1.2cm]{figs/logo/imta.png}\hfill%
	\includegraphics[height=1.2cm]{figs/logo/ls2n.png}\hfill%
	\includegraphics[height=1.5cm]{figs/discovery-sticker.png}\hfill%
 }

% ********************************************************************
% Setup, finetuning, and useful commands
% ********************************************************************
\newcommand{\ie}{i.\,e.}
\newcommand{\Ie}{I.\,e.}
\newcommand{\eg}{e.\,g.}
\newcommand{\Eg}{E.\,g.}

% ************************************************************************
% Loading some handy packages
% ************************************************************************
%% \PassOptionsToPackage{utf8}{inputenc}
%%  \usepackage{inputenc}
\PassOptionsToPackage{english}{babel}
\usepackage{babel}

% ********************************************************************
% General useful packages
% ********************************************************************
\PassOptionsToPackage{T1}{fontenc} % T2A for cyrillics
	\usepackage{fontenc}
\usepackage{xspace} % to get the spacing after macros right

% ************************************************************************
% code / listings
% ************************************************************************
\usepackage{listings}

\lstdefinelanguage{proverif-pv}{
classoffset=0,
morekeywords={%
  among,channel,choice,clauses,const,def,%
  elimtrue,else,equation,equivalence,event,%
  expand,fail,forall,free,fun,get,if,in,%
  inj-event,insert,let,letfun,new,noninterf,%
  not,nounif,or,otherwise,out,param,phase,%
  pred,proba,process,proof,putbegin,query,%
  reduc,set,suchthat,table,then,type,%
  weaksecret,yield},
otherkeywords={||,&&,->,<->,<=>,==>},
emph={gather},
sensitive=true,
morecomment=[n]{(*}{*)},
morestring=[b]",
%rulesepcolor=\color{solarized@base03},
numbers=left,
classoffset=1,
morekeywords={private,data,typeConverter,%
  memberOptim,decompData,decompDataSelect,%
  block,attacker,mess},
keywordstyle=\bfseries\color{solarized@base00},
classoffset=0
}[keywords,comments,strings]


\lstdefinelanguage{idris}[]{Haskell}{
  morekeywords={total,auto,using,Refl,interface,%
  implementation},
  deletekeywords={intersect,List,Double,Bool,%
  String,nub,fst,snd,not,Either,Maybe,Just,Nothing,%
  either,index,last,id,Eq},
  otherkeywords={=>,>>=},
  morecomment=[l]{|||},
  literate=
    {[]}{{\bot}}2 {->}{→}1 {++}{{+\!\!+}}1
    {<-}{{←}}1
    ,
  mathescape=true,% /!\ takes care of '\' inside a code
}

\lstdefinelanguage{javascript}{
  keywords={typeof, new, true, false, catch, function, return, null,
  catch, switch, var, if, in, while, do, else, case, break},
  ndkeywords={class, export, boolean, throw, implements, import,
  this},
  sensitive=false,
  comment=[l]{//},
  morecomment=[s]{/*}{*/},
  morestring=[b]',
  morestring=[b]"
}

\lstset{
basicstyle=\ttfamily\small,
numbers=none,%left,%
numberstyle=\tiny\color{solarized@base01},
keywordstyle=\bfseries\color{solarized@orange},
stringstyle=\color{solarized@cyan},
%% identifierstyle=\color{solarized@base03},
commentstyle=\color{solarized@base01},
%% emphstyle=\color{solarized@orange},
emphstyle=\bfseries\color{solarized@orange},
fontadjust=true,%
keepspaces=true,%
basewidth={0.58em, 0.53em},%
breaklines=true,
tabsize=2,%
%framerule=0.3pt,%
numbersep=1.5mm,%
upquote=true, % straight quote instead of quotation mark
extendedchars=true,
aboveskip=0pt,
belowskip=0pt,
%% mathescape=true,% /!\ this put a mess in the sh code
%% escapechar=★,%
%% firstnumber=last,
escapeinside={`}{'} % CR as escape char
}


% ************************************************************************
% 6. PDFLaTeX, hyperreferences and citation backreferences
% ************************************************************************
% ********************************************************************
% Using PDFLaTeX
% ********************************************************************
%% I got an error when having a pound `#' symbol in the first argument
%% of an `\href' inside a footnote. The footnote tries to interpret
%% it. The bigfoot package fixes the problem.
%% http://tex.stackexchange.com/a/12860
\PassOptionsToPackage{pdftex,hyperfootnotes=false,pdfpagelabels}{hyperref}
	\usepackage{hyperref} % backref linktocpage pagebackref
\pdfcompresslevel=9
% \pdfadjustspacing=1

% ********************************************************************
% Hyperreferences
% ********************************************************************
\hypersetup{%
    %draft,	% = no hyperlinking at all (useful in b/w printouts)
    colorlinks=false, linktocpage=true, pdfstartpage=3, pdfstartview=FitV,%
    % uncomment the following line if you want to have black links (e.g., for printing)
    %colorlinks=false, linktocpage=false, pdfborder={0 0 0}, pdfstartpage=3, pdfstartview=FitV,%
    breaklinks=true, pdfpagemode=UseNone, pageanchor=true, pdfpagemode=UseOutlines,%
    plainpages=false, bookmarksnumbered, bookmarksopen=true, bookmarksopenlevel=1,%
    hypertexnames=true, pdfhighlight=/O,%nesting=true,%frenchlinks,%
    urlcolor=revealjs1, citecolor=webgreen, %pagecolor=RoyalBlue,%
    %urlcolor=Black, linkcolor=Black, citecolor=Black, %pagecolor=Black,%
    pdftitle={\myTitle},%
    pdfauthor={\textcopyright\ \myName, \myDepartment},%
    pdfsubject={},%
    pdfkeywords={},%
    pdfcreator={pdfLaTeX},%
    pdfproducer={LaTeX with hyperref and classicthesis}%
}


% ********************************************************************
% Last, but not least...
% ********************************************************************
\usepackage{ulem} %to strike out text
\renewenvironment{quote}
{\vspace{.25em}\itshape\list{}{\rightmargin\leftmargin}%
\item\relax\textquotedblleft\ignorespaces}
{\unskip\unskip\textquotedblright\endlist\vspace{.25em}}
% ************************************************************************

% ********************************************************************
% Using different fonts
% ********************************************************************

%% Crimson: https://fontlibrary.org/en/font/crimson
\usepackage[no-math]{fontspec}
%% Scale to math text size
\defaultfontfeatures{Scale=1.05 }
%\newfontfamily\Fonti[
%    Path              = /Users/msimonin/Library/Fonts/,
%    Extension         = .otf,
%    Ligatures         = TeX,
%    UprightFont       = *-Roman,
%    BoldFont          = *-Bold,
%    ItalicFont        = *-Italic,
%    BoldItalicFont    = *-BoldItalic
%]{Crimson}
%\newfontfamily\Fontii[
%    Path              = /Users/msimonin/Library/Fonts/,
%    Extension         = .otf,
%    Ligatures         = TeX,
%    WordSpace         = 1.3,
%    ItalicFont        = Vollkorn-MediumItalic,
%    BoldFont          = Vollkorn-Semibold,
%    BoldItalicFont    = Vollkorn-SemiboldItalic
%]{Vollkorn-Medium}
%\newfontfamily\FontiiD[
%    Path              = /Users/msimonin/Library/Fonts/,
%    Extension         = .otf,
%    Ligatures         = TeX,
%    UprightFont       = *-Regular,
%    BoldFont          = *-Bold,
%    ItalicFont        = *-Italic,
%    BoldItalicFont    = *-BoldItalic
%]{Vollkorn}
%\newfontfamily\Fontiii{Lato}
%
%\setmainfont[
%    Path              = /Users/msimonin/Library/Fonts/,
%    Extension         = .otf,
%    Ligatures         = TeX,
%    UprightFont       = *-Roman,
%    BoldFont          = *-Bold,
%    ItalicFont        = *-Italic,
%    BoldItalicFont    = *-BoldItalic
%]{Crimson}
%
%\setsansfont{Lato}
%
%\setmonofont[
%    Path              = /Users/msimonin/Library/Fonts/,
%    Extension         = .otf,
%    Scale             = MatchLowercase,
%    Ligatures         = TeX,
%    UprightFont       = *-Regular,
%    BoldFont          = *-Medium
%]{FiraMono}%% \linespread{1.05} % a bit more for Palatino
